\scanpage{41}
\mspage{40}
telle famille, $\exists i$ t.q. $F_i=F$, est \inserted{donc} pareil dans $\unsure{\mathcal F|\mathscr F}$ et $\mathcal F'$. )

Supposons que $\mathcal F$ satisfasse At 2 :
\[
\begin{aligned}
\forall F\in\mathcal F,\ \text{on a}\qquad
 F&=\mathop{\mathrm{Sup}}\limits_{X\in\deleted{\unsure{F|\mathcal F}}\widetilde F\unsure{.|,}}X
 &\qquad \text{où}\quad \widetilde F&=\{X\in M\mid X\leq F\}\\
 &&&=\{X\in\mathcal F,\ X\triangleleft F\}.
\end{aligned}
\]
Alors $\mathcal F'$ y satisfait aussi, car si $F\in\mathcal F'$, $\widetilde F$ au sens de $\mathcal F$ ou de $\mathcal F'$ c'est pareil, et le sup dans $\mathcal F$ est un sup dans $\mathcal F'$.

\uline{Définition} \quad La partie fermée $\mathcal F'$ \deleted{\unsure{ou de|est de}} de $\mathcal F$ est dite \uline{spacieuse} si elle est stable par $\mathrm{Sup}$ quelconques (de familles majorées \uline{dans $\mathcal F$}).
% slanted
\margin{On dira qu'elle est \uline{finiment} \uline{spacieuse} si elle est stable par $\mathrm{Sup}$ finis.}
% slanted
\margin{\unsure{Remords|Remonde} : il vaut mieux réserver le mot ``spacieux'' à la notion la plus modeste des deux, et appeler l'autre ``infiniment spacieux'' ou ``très spacieux''.}

\uline{Exemple} \quad Soit $\mathcal F$ \inserted{$=\unsure{\mathcal P|\mathfrak P}(I)$} l'ens des parties d'un ens. $I$, \deleted{Alors $\mathcal F'$} avec la relation d'inclusion. Satisfait At 1, At 2, At 3, \deleted{\unsure{les|la}} $M=\{\{i\}\mid i\in I\}$ (avec \unsure{ordre|l'ordre} discret), \inserted{(si $I$ est fini)}
$\deleted{\unsure{\mathcal F|I}}$Se donner une $\mathcal F'$ \inserted{\deleted{fer\unsure{mée|mé}} \uline{\unsure{ample|sp}}} revient au $=$ que de se donner une structure de \deleted{\unsure{pr|s}} \deleted{\unsure{quelque|quels|quasi}} simpliciale sur \inserted{une partie $\unsure{I'|I}\deleted{\unsure{\subset|=}}=\mathcal F'\cap I$ de} $I$, du moins si $I$ est fini. (Sinon, c'est vrai à condition de se restreindre aux $\mathcal F'$ \unsure{formés|fermés} d'ens. finis \ldots).

Dire que $\mathcal F'$ est spacieux, signifie que \deleted{\unsure{tout|toute}} réunion de simplexes est une simplexe -- se donner un tel $\mathcal F'$, revient donc à se donner une partie $I'$ de $I$, \deleted{$\mathcal F'$ est} et $\deleted{\unsure{F|\mathcal J}}\mathcal F'=\mathcal P(I')$ -- c'est pas du \deleted{tout} pareil !
